% generated by roam.py handbook export -- не редактировать, править заметки
% 20260922175350-связь_спектральнои_нормы_гессиана_и_липшицевости_градиента.org (body 083af92f09d9)

\begin{theorem}[{Связь спектральной нормы гессиана и липшицевости градиента}]\label[theorem]{note:4B6D53BA-E599-4DEA-A557-550BD32042BA}
  Пусть $f: E \rightarrow \R$, $E$~--- \hyperref[note:316B5BFA-8B9F-4267-9086-D6593AAB5F9B]{выпуклое} и $f \in C^2$ \\
  И $\exists L:\ \forall x\ \| \nabla^2 f(x) \| \le L$, где $\| \cdot \|$~--- спектральная норма (\hyperref[note:592AD816-1D80-48E6-8399-8524FBC1E8F3]{гессиан}) \\
  Тогда $f \in C^{1,1}_L$
\end{theorem}

\begin{remark}
  Верно и обратное~--- \hyperref[note:DAA0DAE4-4B48-4A2F-9459-54790B6B129F]{оценка нормы гессиана сверху}, так что для $f \in C^2$: $\nabla f$~--- $L$-липшицев $\Leftrightarrow \| \nabla^2 f(x) \| \le L\ \forall x$.
\end{remark}

\begin{remark}
  $\nabla^2 f(x)$ симметрична, поэтому по \hyperref[note:3BE0BBCE-AAB0-4CB4-8FA6-45277C0889CF]{спектральной теореме} $\| \nabla^2 f(x) \| = \max_i \abs{\lambda_i}$, и условие равносильно $-L I \le \nabla^2 f(x) \le L I$. Для $\mu$-сильно выпуклой $f$ получаем $\mu I \le \nabla^2 f(x) \le L I$, откуда снова видно $\mu \le L$.
\end{remark}

\begin{proof}
  Вспомним \hyperref[note:5821D7FE-1F9F-4578-BA2B-D60FDB0C69D3]{теорему Лагранжа для векторнозначных функций}:
  Пусть $g: E \rightarrow \R^{n}$,  $E$~--- выпукло и $g$~--- дифференцируема на $E$ \\
  Тогда $\forall a,b \in E\ \exists c \in [a;b]:\ \| g(b) - g(a) \| \le \| g^{\prime}(c) \| \cdot \| b -a \|$

  Рассмотрим $g(x) := \nabla f(x)$ и получим
  \[
  \| \nabla f(a) - \nabla f(b) \| \le \| \nabla^2 f(c) \| \cdot \| b - a \|  \le  L \| b - a \|
  \]
  Получили определение \hyperref[note:7DD1C7AB-1A49-4675-9F21-5A16C717F670]{липшицевости градиента}
\end{proof}
